\section{A variational GLLVM re-estimation: model, inference, and comparison to NUTS}
\label{ann:variational}

The atlas of the main text is estimated by Hamiltonian Monte Carlo (the No-U-Turn sampler,
NUTS), which is audit-grade but costs hours of multi-chain sampling. Here we re-estimate the
\emph{same} measurement model by stochastic variational inference (SVI), giving a
${\sim}100\times$ faster operational engine that we calibrate against---rather than substitute
for---the NUTS fit. We frame this as a generalized linear latent-variable model
(GLLVM)~\citep{hui2017gllvm} fit variationally~\citep{blei2017vi,hoffman2013svi}; a fuller
development of the variational engine (amortized scoring, richer posteriors, calibrated
uncertainty) is left to future work. This section describes the model, the inference, the
conceptual differences from NUTS, and the head-to-head comparison.

\subsection{Generative model}

The generative model is the atlas of Annex~\ref{ann:measurement}, unchanged. For patient $i$
with factor vector $f_i=(G_i,D_{i1},\dots,D_{iK})^{\top}$ and item $j$,
\begin{align}
  f_i &\sim \mathcal N(\mathbf 0,\Phi), &&\Phi\ \text{bifactor-faithful (\Gfac{} and the pinned substance axis orthogonal),} \label{eqH:prior}\\
  \eta_{ij} &= \alpha_j + \lambda_j^{\top} f_i, &&\text{linear decoder; home loadings positive (softplus), forbidden cells exact zero,}\label{eqH:decoder}\\
  x_{ij} &\sim F_j(\eta_{ij},\theta_j), &&\text{per-item likelihood,}\label{eqH:lik}
\end{align}
with the same ontology prior matrix as the folded operational map (immunometabolic merge; the
three earned cross-loadings freed; everything else hard-zero). $F_j$ is Gaussian on the
rank-INT (Gaussian-copula) scale for continuous and high-cardinality items, and native
Bernoulli-logit / ordered-logistic / negative-binomial for binary, low-cardinality ordinal and
count items---the same tiered likelihood as the main fit. The log-likelihood is evaluated over
\emph{observed cells only} (a missingness mask removes absent cells; no imputation). One
implementation difference from the NUTS engine: rather than marginalizing the continuous block
analytically (the Woodbury identity of Annex~\ref{ann:measurement}), the variational engine
retains an explicit per-patient latent for \emph{every} factor and attaches every likelihood
directly to $f_i$, so it carries the full $\approx142$-indicator block with no
continuous/explicit split.

\subsection{Variational inference}

The exact posterior $p(\{f_i\},\Theta\mid X_{\mathrm{obs}})$ is intractable, so SVI replaces it
with a tractable family and \emph{optimizes} rather than samples. Each patient receives an
independent variational posterior over its coordinate,
\begin{equation}
  q_i(f_i)=\mathcal N(\mu_i,\Sigma_i),\qquad
  \Sigma_i=\operatorname{diag}(s_i^2)\ \text{(mean-field default)},
  \label{eqH:q}
\end{equation}
with optional low-rank-plus-diagonal $\Sigma_i=\operatorname{diag}(s_i^2)+U_iU_i^{\top}$ or full
$K\times K$ covariance. Inference minimizes the negative evidence lower bound (ELBO),
\begin{equation}
  \mathcal J \;=\; \underbrace{\mathbb E_{q}\!\left[-\log p(X_{\mathrm{obs}}\mid f,\Theta)\right]}_{\text{observed-cell reconstruction}}
  \;+\; \sum_{i=1}^{N}\operatorname{KL}\!\left[\,q_i(f_i)\,\big\|\,\mathcal N(\mathbf 0,\Phi)\,\right]
  \;+\; \Omega(\Lambda) + \Omega(\Phi),
  \label{eqH:elbo}
\end{equation}
where the reparameterization trick~\citep{kingma2014vae} makes the latent draws differentiable,
the KL term regularizes each coordinate toward the bifactor prior \eqref{eqH:prior}, and
$\Omega(\cdot)$ are the ontology penalties---the soft loading prior written as a per-cell
$\tfrac12\sum_{\text{free}}((\lambda-m)/v)^2$ with $(m,v)$ read from the prior matrix, plus a
weak off-diagonal penalty on $\Phi$. The objective is optimized by minibatch AdamW (per-patient
terms scaled by $N/B$; global penalties added once), with the small $K\times K$ Cholesky/KL on
CPU in double precision. The loadings, factor correlations, residual scales and cutpoints
($\Lambda,\Phi,\sigma,\text{cutpoints}$) are \emph{maximum-a-posteriori} point estimates;
loading credible bands are obtained from a seed ensemble, with NUTS retained as the
uncertainty authority. A two-rung warm start (continuous backbone $\to$ full mixed map)
protects the thin factors from ELBO-flat collapse. The full fit takes ${\approx}4.5$~minutes on
a workstation CPU (PyTorch~\citep{paszke2019pytorch}), against the NUTS fit's hours.

\subsection{NUTS versus variational inference}

The two engines optimize the same model but differ in what they compute and what they cost
(Table~\ref{tab:nuts_vs_vi}). NUTS draws asymptotically exact samples from the full joint
posterior, delivering calibrated uncertainty for every quantity; SVI returns the best
approximation \emph{within a chosen family}, which is fast but biased by the family. The
load-bearing consequence here is specific and predictable: a mean-field (diagonal) $q_i$ cannot
represent posterior covariance \emph{between} factors, so to fit the data it shrinks the prior
inter-factor correlations the likelihood would otherwise support. The result is faithful factor
loadings and patient coordinates---which are governed by the posterior \emph{means}---but
systematically attenuated inter-factor correlations $\Phi$ and underestimated coordinate
standard deviations, both governed by the posterior \emph{covariance}. This is why we keep NUTS
as the authority for $\Phi$ magnitudes and for any credible interval, and use the variational
engine for the map, the coordinates, fast reruns and synthetic-cohort generation.

\begin{table}[ht]\centering\small
\caption{NUTS versus variational inference for the FACE-ATLAS measurement model.}
\label{tab:nuts_vs_vi}
\begin{tabular}{@{}p{3.2cm}p{5.4cm}p{5.4cm}@{}}
\toprule
 & \textbf{NUTS (main fit)} & \textbf{Variational GLLVM} \\
\midrule
Inference & MCMC: asymptotically exact samples from the full joint posterior & Optimization: best approximation within a chosen posterior family \\
Latent treatment & Continuous block marginalized (Woodbury); latents only for discrete items & Explicit per-patient latent for every factor; all $\approx142$ items \\
Global parameters & Full posterior (loadings, $\Phi$, residuals) with credible intervals & MAP point estimates; loading bands from a seed ensemble \\
Coordinate uncertainty & Calibrated posterior SD per axis & Relative reliability signal (mean-field underestimates SD) \\
Inter-factor $\Phi$ & Posterior magnitudes with uncertainty & Structure correct; off-diagonals attenuated (mean-field bias) \\
Cost & Hours (multi-chain) & ${\approx}4.5$ min CPU (${\sim}100\times$ faster) \\
\bottomrule
\end{tabular}
\end{table}

\subsection{Comparison results}

On the eight-factor operational map, the variational engine reproduces the NUTS fit on its two
load-bearing objects---the measurement map and the patient coordinates---and matches it as a
generative model.

\paragraph{The measurement map.} Loading columns are congruent on all eight factors (Tucker
$\varphi$: \Gfac{} $0.957$, every specific axis $0.96$--$0.999$; Table~\ref{tab:vi_loadings}),
and the full item$\times$factor loading scatter is $r=0.993$ over $208$ cells with no sign flips
(Figure~\ref{fig:vi_loadings}). \Gfac{} is the hardest axis---it carries every item's small
bifactor loading---and is the only one near the major-axis bar; a richer posterior raises it to
$0.974$. \emph{It is the same map.}

\paragraph{The patient coordinates.} Per-patient coordinates correlate with the NUTS
coordinates at $r\ge0.90$ on all eight axes (Table~\ref{tab:vi_coords}: \Gfac{} $0.990$,
immunometabolic $0.998$, cognition $0.992$, down to substance $0.908$, the orthogonal,
instrument-poor axis). \emph{Patients land in the same place.}

\begin{table}[ht]\centering\small
\begin{minipage}[t]{0.48\textwidth}\centering
\caption{Loading congruence (Tucker $\varphi$) vs NUTS.}
\label{tab:vi_loadings}
\begin{tabular}{@{}lrr@{}}
\toprule
Factor & cells & Tucker $\varphi$ \\
\midrule
\Gfac{} (overall severity) & 106 & 0.957 \\
cognition & 16 & 0.998 \\
immunometabolic & 37 & 0.998 \\
sleep & 12 & 0.999 \\
suicidality & 12 & 0.998 \\
developmental risk & 21 & 0.998 \\
mania/activation & 2 & 0.984 \\
substance & 2 & 0.963 \\
\bottomrule
\end{tabular}
\end{minipage}\hfill
\begin{minipage}[t]{0.48\textwidth}\centering
\caption{Coordinate agreement (Pearson $r$) vs NUTS.}
\label{tab:vi_coords}
\begin{tabular}{@{}lrr@{}}
\toprule
Factor & $n$ (well) & $r$ \\
\midrule
\Gfac{} (overall severity) & 8641 & 0.990 \\
cognition & 6451 & 0.992 \\
immunometabolic & 8825 & 0.998 \\
sleep & 7522 & 0.968 \\
suicidality & 8285 & 0.984 \\
developmental risk & 8767 & 0.986 \\
mania/activation & 6152 & 0.976 \\
substance & 4249 & 0.908 \\
\bottomrule
\end{tabular}
\end{minipage}
\end{table}

\begin{figure}[ht]
  \centering
  \includegraphics[width=0.6\linewidth]{gllvm_oop/vi_vs_nuts_loadings.png}
  \caption{\textbf{The variational map is the NUTS map.} Item$\times$factor loadings from the
  variational GLLVM against the NUTS copula fit ($n=208$ cells, $r=0.993$), coloured by cell
  type. Large home and \Gfac{}-anchor loadings align on the diagonal; cross, window and small
  bifactor-\Gfac{} cells cluster tightly near zero.}
  \label{fig:vi_loadings}
\end{figure}

\paragraph{Generative fidelity.} Drawing $9{,}013$ synthetic patients from the fitted map and
comparing to the observed data, the marginals match (Kolmogorov--Smirnov median $0.04$;
per-variable mean fidelity $r=1.000$, SD fidelity $r=0.998$) and so does the joint covariance
(correlation SRMR $0.078$, matching the NUTS fit's $\approx0.07$)---the eight factors reproduce
the data's correlation structure, not merely its margins.

\paragraph{The one gap: inter-factor $\Phi$.} The correlation structure is reproduced---\Gfac{}
and substance rows are exactly zero in both engines, and the
sleep/suicidality/developmental-risk/mania block is positively correlated in both
(Figure~\ref{fig:vi_phi})---but the off-diagonals are systematically attenuated: correlated-block
mean $|\phi|$ is $0.083$ (variational) versus $0.105$ (NUTS), a ${\sim}21\%$ shrinkage
concentrated in the small negative cognition correlations (the substantive
immunometabolic$\leftrightarrow$cognition coupling is preserved). Climbing the full
posterior-richness ladder---mean-field $\to$ low-rank $\to$ full $K\times K$ covariance---reduces
the shrinkage monotonically ($21\%\to20\%\to18\%$) but does not close it, establishing the
residual as a \emph{structural bias of variational inference on the inter-factor correlation
hyperparameter} rather than a tunable defect (an exploratory arm removing the $\Phi$ penalty and
loosening gradient clipping left $\Phi$ unchanged). This is precisely the quantity for which
NUTS is retained as the authority; the variational $\Phi$ should be read as directionally
correct but not as a magnitude.

\begin{figure}[ht]
  \centering
  \includegraphics[width=\linewidth]{gllvm_oop/vi_vs_nuts_phi.png}
  \caption{\textbf{Inter-factor correlation $\Phi$: structure reproduced, magnitudes attenuated.}
  Variational $\Phi$ (left), NUTS $\Phi$ (middle), and their difference (right). The
  block-diagonal orthogonality (\Gfac{}, substance) is exact in both and the correlated specific
  block is positive in both; the difference panel localizes the ${\sim}21\%$ attenuation to the
  small cognition cross-correlations.}
  \label{fig:vi_phi}
\end{figure}

\paragraph{Summary.} The variational GLLVM is a calibrated, ${\sim}100\times$-faster stand-in for
the NUTS measurement map and patient coordinates (loadings Tucker $8/8$, scatter $r=0.993$;
coordinates $r\ge0.90$ on $8/8$ axes) and a faithful generative model of the cohort, with the
inter-factor correlation $\Phi$ as the one structurally attenuated quantity for which NUTS
remains the authority. It is suited to fast reruns, sensitivity sweeps, synthetic-cohort
generation, and as an independent-estimator robustness check on the map; developing it into a
fully calibrated Bayesian engine is future work.
